\section{Discussion, Limitations, and Outlook}
\label{sec:discussion}

This work reframes anisotropy in human language representations as a dual
phenomenon.
On one hand, the presence of a dominant spectral spike poses significant
challenges for optimization and quantization.
On the other hand, the coherence and shared nature of this spike make it
amenable to simple, efficient suppression through averaging.
By grounding this perspective in empirical analysis, a minimal structural
model, and formal guarantees, we show that averaging can fill the role of
spectral spike suppression in a principled and hardware-friendly manner.

\subsection{When Averaging Works}

The effectiveness of averaging relies on structural properties that are
ubiquitous in large language model training, but not universal.
In particular, averaging is most effective when:
\begin{itemize}
  \item the dominant spike corresponds to shared linguistic structure that is
  coherent across tokens;
  \item semantic components are diverse and weakly correlated across sequences
  or sufficiently distant contexts;
  \item batches contain a sufficient number of semantically independent tokens
  so that semantic residuals cancel under averaging.
\end{itemize}

These conditions are empirically validated for standard language model training
regimes, but may not hold in settings with highly correlated batches, extremely
small batch sizes, or non-linguistic data distributions.
In such cases, more aggressive spectral methods may still be required.

\subsection{Conservatism versus Aggressiveness}

A key theme throughout the paper is that averaging is a \emph{conservative}
operation.
Unlike SVD-based spike suppression, which removes an entire low-rank subspace,
averaging removes only a single coherent direction aligned with the batch mean.
This conservatism is both a strength and a limitation.

On the positive side, it minimizes collateral damage to semantic structure and
avoids amplifying noise in the tail.
On the negative side, it cannot eliminate all spike energy when the spike
subspace has dimension greater than one.
As a result, averaging should be viewed as a safe first-order approximation,
rather than a complete replacement for spectral methods in all settings.

\subsection{Relation to Existing Spectral Methods}

Our analysis clarifies why explicit spectral methods such as Metis and
spike-aware optimizers are effective: they target the same coherent structure
identified here.
The contribution of this work is not to dispute the value of such methods, but
to show that much of their benefit can be achieved by exploiting statistical
regularities of language data through simple operations.

From this perspective, averaging can be seen as a baseline spectral
approximation.
It may be combined with more refined methods, used as a warm-up or
preconditioning step, or applied selectively in regimes where full spectral
analysis is impractical.

\subsection{Broader Implications}

The curse–blessing duality highlighted in this work suggests a broader principle
for large-scale learning systems:
pathological geometric structures induced by data may also encode exploitable
regularities.
Understanding the geometry of data distributions—not just model architectures—
can open the door to algorithmic simplifications that are both theoretically
sound and practically impactful.

While this paper focuses on language models, similar anisotropic structures have
been observed in vision, multimodal models, and scientific data.
Exploring whether analogous averaging-based approximations apply in these
domains is an interesting direction for future work.

\subsection{Conclusion}

Anisotropy in human language corpora is unavoidable.
It degrades training efficiency, destabilizes quantization, and complicates
optimization.
At the same time, its coherence enables simple and efficient remedies.

This work shows that averaging, one of the simplest operations available in
large-scale systems, can be elevated from a heuristic to a principled spectral
approximation.
By aligning algorithmic design with the geometry of language data, we obtain a
method that is efficient, interpretable, and broadly applicable.
We hope this perspective encourages further exploration of data-aware
approaches to large-scale model training.